% ============================================================================
% Курсов проект по дисциплина „Информационни технологии за бизнес мениджмънт“
% ОКС „Магистър“, ФКСТ, ТУ-София, I семестър.
%
% Тема: система за подпомагане на решенията при избор и подреждане на портфейл
%       от ИТ проекти при бюджетни и ресурсни ограничения.
%
% Компилация:  latexmk -pdf Main.tex        (pdfLaTeX, НЕ XeLaTeX)
% Правила:     форматът следва изискванията на ФКСТ, преведени в preamble.tex
% Форматиране: preamble.tex
%
% Този файл е СКЕЛЕТЪТ: структурата и редът на разделите. Съдържанието живее
% в chapters/. Незапълнените места са маркирани с \TODO{...} и се намират с
%     grep -rn 'TODO' chapters/ Main.tex references.bib
% ============================================================================
\documentclass[12pt]{article}

% ============================================================================
% preamble.tex: форматиране по изискванията на ФКСТ.
%
% Всяка секция тук съответства на правило от REQUIREMENTS.md (Част А) и на
% реда му в таблицата на превода (Част Б). Не пипай стойностите без да
% провериш срещу REQUIREMENTS.md, те са нормативни, не са вкус.
%
% Двигател: pdfLaTeX (T2A + tempora). НЕ XeLaTeX, вж. REQUIREMENTS.md Б.0.
% ============================================================================

\usepackage{cmap}
\usepackage[T2A]{fontenc}
\usepackage[utf8]{inputenc}
\usepackage[bulgarian]{babel}

% --- Типография (А.4) ---------------------------------------------------
% tempora = Times-метрично съвместим шрифт с пълна кирилица под T2A.
\usepackage{tempora}
\usepackage[top=2.5cm, bottom=2.5cm, left=3.0cm, right=2.0cm, a4paper]{geometry}
\usepackage{setspace}
\setstretch{1.5}                     % „междуредие 1.5 lines“, калибруем, вж. Б.2
% „без пренасяния на срички“ (А.4), чрез ЗАБРАНИТЕЛНИ НАКАЗАНИЯ, не чрез
% \usepackage[none]{hyphenat}. Разликата е съществена: `hyphenat[none]` задава
% \hyphenchar=-1 и така изключва самия механизъм за прекъсване, от който зависи
% `breaklines` на listings: дългите редове в листингите тогава НЕ се пренасят и
% излизат извън полето (1227 препълвания дойдоха оттам). Наказанията постигат същия
% резултат за текста (TeX никога не пренася), но оставят механизма достъпен.
\usepackage[expansion=false]{microtype}
\DisableLigatures{encoding=*,family=tt*}
\hyphenpenalty=10000
\exhyphenpenalty=10000
\emergencystretch=1.5em              % цената на justified-без-преноси, вж. Б.4
\sloppy
\usepackage{indentfirst}             % отстъп и на ПЪРВИЯ абзац след заглавие, вж. Б.3
\setlength{\parindent}{1.25cm}
\setlength{\parskip}{0pt}

% --- Номерация на страниците: долу вдясно (А.4) -------------------------
\usepackage{fancyhdr}
\pagestyle{fancy}
\fancyhf{}
\renewcommand{\headrulewidth}{0pt}
\fancyfoot[R]{\thepage}

% --- Бележки под линия: 10 pt, single, justified (А.4) -------------------
\usepackage[hang,flushmargin]{footmisc}
\usepackage{ragged2e}
\renewcommand{\footnotelayout}{\fontsize{10}{12}\selectfont\setstretch{1.0}\justifying}

% --- Заглавия на раздели / параграфи / подпараграфи (А.7) ---------------
\usepackage{titlesec}
\newcommand{\sectionbreak}{\clearpage}   % „всеки раздел започва на нова страница“

\renewcommand{\thesection}{\Roman{section}}
\renewcommand{\thesubsection}{\arabic{section}.\arabic{subsection}}
\renewcommand{\thesubsubsection}{\arabic{section}.\arabic{subsection}.\arabic{subsubsection}}

% Раздел: римски, ALL CAPS, центриран, bold 14 pt
\titleformat{\section}[block]
  {\fontsize{14}{16.8}\bfseries\centering}{\thesection. }{0pt}{\MakeUppercase}

% Параграф: арабски, bold 14 pt, отстъп 1.25 cm
\titleformat{\subsection}[block]
  {\fontsize{14}{16.8}\bfseries}{\thesubsection. }{0pt}{}
\titlespacing*{\subsection}{1.25cm}{*3}{*2}

% Подпараграф: bold 13 pt, отстъп 2.0 cm
\titleformat{\subsubsection}[block]
  {\fontsize{13}{15.6}\bfseries}{\thesubsubsection. }{0pt}{}
\titlespacing*{\subsubsection}{2.0cm}{*2}{*1}

% --- Съдържание: разделите с ГЛАВНИ букви, 12 pt, bold (А.6) -------------
\usepackage{etoolbox}
\usepackage[titles]{tocloft}
% NB: кръпката ТРЯБВА да е СЛЕД tocloft, пакетът предефинира \l@section и
% иначе я изтрива (тихо: няма грешка, просто разделите излизат с малки букви).
\makeatletter
\patchcmd{\l@section}{#1}{\MakeUppercase{#1}}{}{}
\makeatother
\renewcommand{\cftsecleader}{\cftdotfill{\cftdotsep}}
\renewcommand{\cftsubsecleader}{\cftdotfill{\cftdotsep}}
\renewcommand{\cftsubsubsecleader}{\cftdotfill{\cftdotsep}}
\renewcommand{\cftsecfont}{\fontsize{12}{14}\bfseries}
\renewcommand{\cftsecpagefont}{\fontsize{12}{14}\bfseries}
\renewcommand{\cftsecaftersnum}{.\enspace}
\settowidth{\cftsecnumwidth}{VIII.\enspace}
\renewcommand{\cftsubsecaftersnum}{.}
\renewcommand{\cftsubsubsecaftersnum}{.}

% --- Надписи: таблици ОТГОРЕ, фигури ОТДОЛУ (А.8, А.9) ------------------
\usepackage{caption}
\usepackage{subcaption}
\DeclareCaptionLabelFormat{tablelabel}{Таблица #2}
\captionsetup[table]{labelformat=tablelabel, labelsep=period, position=above,
                     font={bf,small}, justification=centering, singlelinecheck=true}
\DeclareCaptionLabelFormat{figurelabel}{Фигура #2}
\captionsetup[figure]{labelformat=figurelabel, labelsep=period, position=below,
                      font={bf,small}, justification=centering, singlelinecheck=true}

% Препратки в текста: „(Табл. 3)“ / „(Фиг. 7)“, задължителни по А.8/А.9.
\newcommand{\tabref}[1]{(Табл.~\ref{#1})}
\newcommand{\figref}[1]{(Фиг.~\ref{#1})}

% --- Математика ---------------------------------------------------------
\usepackage{amsmath,amssymb,amsfonts}
\usepackage{amsthm}

% --- Графики ------------------------------------------------------------
\usepackage{graphicx}
\usepackage{float}
\usepackage{tikz}
\usetikzlibrary{decorations.pathreplacing,arrows.meta,positioning,shapes.geometric,fit,calc}
\usepackage{adjustbox}
\usepackage{pdfpages}     % за подвързаните задание и декларация, вж. Б.5

% --- Таблици ------------------------------------------------------------
\usepackage{array}
% Ляво подравнена колона с фиксирана ширина. Тесните `p{}` колони с ДВУСТРАННО
% подравняване (наследено от основния текст) разкъсват редовете с огромни
% междудумни интервали: правилото за justified от А.4 се отнася за ОСНОВНИЯ текст,
% не за клетки от 3 cm. Всички таблици с текстови клетки ползват L.
\newcolumntype{L}[1]{>{\RaggedRight\arraybackslash}p{#1}}
\usepackage{booktabs}
\usepackage[table]{xcolor}% \cellcolor за топлинната карта
\usepackage{longtable}    % „продължение“ при пренасяне на таблица (А.8)
\usepackage{multirow}
\usepackage{tabularx}

% --- Помощни ------------------------------------------------------------
\usepackage{hyperref}
\usepackage{xurl}
\usepackage{enumitem}
\usepackage{textcomp}
\usepackage{listings}
\urlstyle{same}
% КРИТИЧНО: `breaklines` на listings вмъква ДИСКРЕЦИОННИ прекъсвания, а глобалните
% \hyphenpenalty/\exhyphenpenalty=10000 (правилото „без пренасяния“, А.4) ги правят
% безкрайно скъпи, тогава дългите редове в листингите НЕ се пренасят и излизат извън
% полето. Проверено изолирано: с наказанията 1 препълване, без тях 0. Наказанията се
% нулират САМО вътре в листинг, така че текстът си остава без пренасяния.
\makeatletter
\lst@AddToHook{PreInit}{\hyphenpenalty=50\exhyphenpenalty=50}
\makeatother
\renewcommand{\lstlistingname}{Листинг}
\renewcommand{\lstlistlistingname}{Списък на листингите}

% Italic САМО за ключови понятия (А.4): семантичен макрос, за да е ревизуемо.
\newcommand{\term}[1]{\emph{#1}}
% Инлайн код/идентификатори.
\DeclareRobustCommand{\code}[1]{\texttt{#1}}
% Глифи за таблиците с присъди (генерирани от tools/gen_tex.py).
\newcommand{\cmark}{\ensuremath{\checkmark}}
\newcommand{\xmark}{\ensuremath{\times}}

% Числата от борда са генерирани, за да не може изречение да противоречи на таблица.
% Заместващ файл. preamble.tex го зарежда безусловно, затова той трябва да
% съществува, за да тръгне компилацията, докато tools/ още не произвежда
% generated/. Празен е нарочно: не дефинира нито един макрос с число.
% Когато измерванията бъдат направени, този файл се презаписва от инструмента и
% дефинициите тук идват от данните, а не от ръка.

% Заместващ файл, вж. бележката в board_macros.tex. Празен е нарочно.


% Заснемане с тънка черна рамка. Рамката НЕ е украса: интерфейсите със светла тема са
% почти бели и се сливат с бялата страница, така че без нея не се вижда къде свършва
% изображението и започва листът, а размерът на изображението е част от твърдението.
% \fboxsep=0pt прилепва рамката към пиксела; 0.4 pt е най-тънката, която се възпроизвежда
% надеждно при печат.
\newcommand{\capture}[2][0.44\textwidth]{%
  {\setlength{\fboxsep}{0pt}\setlength{\fboxrule}{0.4pt}%
   \fbox{\includegraphics[width=#1]{#2}}}}

% Съпоставка „еталон | система“ за един интерфейс.
%   #1 лента  #2 ключ на интерфейса  #3 тема  #4 описание
\newcommand{\evidencepair}[4]{%
  \begin{figure}[H]\centering
    \begin{tabular}{@{}cc@{}}
      \small\textbf{еталон} & \small\textbf{система} \\
      \capture{evidence/#1/maui/#2_#3.png} &
      \capture{evidence/#1/cpp/#2_#3.png} \\
    \end{tabular}
    \caption{#4}\label{fig:ev-#1-#2-#3}
  \end{figure}}

% Видим маркер за незапълнено място. Търси се с: grep -rn 'TODO' chapters/
\newcommand{\TODO}[1]{\textcolor{red}{\textbf{[TODO: #1]}}}

\tikzset{
    block/.style={draw, rounded corners, fill=blue!5, align=center, font=\small,
        minimum width=3.2cm, minimum height=1.05cm, inner sep=6pt},
    accent/.style={draw, rounded corners, fill=green!8, align=center, font=\small,
        minimum width=3.0cm, minimum height=0.95cm, inner sep=6pt},
    danger/.style={draw, rounded corners, fill=red!6, align=center, font=\small,
        minimum width=3.0cm, minimum height=0.95cm, inner sep=6pt},
    note/.style={draw, dashed, rounded corners, fill=gray!8, align=left, font=\small,
        inner sep=6pt},
    line/.style={-Latex, thick}
}

% --- Единично разстояние в нетекстови среди -----------------------------
\AtBeginEnvironment{tikzpicture}{\setstretch{1.1}}
\AtBeginEnvironment{lstlisting}{\setstretch{1.1}}
\AtBeginEnvironment{verbatim}{\setstretch{1.1}}
\AtBeginEnvironment{tabular}{\setstretch{1.1}}
\AtBeginEnvironment{longtable}{\setstretch{1.1}}

\BeforeBeginEnvironment{tikzpicture}{\begin{adjustbox}{max width=\textwidth}}
\AfterEndEnvironment{tikzpicture}{\end{adjustbox}}

% --- Оцветяване на C++ код ----------------------------------------------
\definecolor{cppblue}{RGB}{0,0,180}
\definecolor{cppgreen}{RGB}{0,128,0}
\definecolor{cppred}{RGB}{180,0,0}
\definecolor{cppgray}{RGB}{128,128,128}

\lstset{
    language=C++,
    basicstyle=\footnotesize\ttfamily,
    % Пренасяне на дълги редове. `breakatwhitespace=false` е СЪЩЕСТВЕНО: коментарите в
    % изходния текст стигат 110 знака без удобно място за прекъсване, и при `true` редът
    % изобщо не се пренася, а излиза извън полето. `linewidth` се задава изрично, защото
    % рамката и левият отстъп иначе не се приспадат от наличната ширина.
    breaklines=true, breakatwhitespace=false, breakautoindent=false, breakindent=0pt,
    linewidth=\dimexpr\linewidth-2.5em\relax,
    columns=flexible, keepspaces=true,
    postbreak=\mbox{\textcolor{cppgray}{\ensuremath{\hookrightarrow}}\space},
    xleftmargin=1.5em, xrightmargin=0.5em,
    keywordstyle=\color{cppblue}\bfseries,
    stringstyle=\color{cppred},
    commentstyle=\color{cppgreen},
    identifierstyle=\color{black},
    showstringspaces=false,
    numbers=left, numberstyle=\tiny\color{cppgray}, numbersep=8pt,
    aboveskip=1.5em, belowskip=1.5em,
    frame=single, framerule=0.5pt, rulecolor=\color{cppgray},
    tabsize=4, captionpos=b,
    morekeywords={concept, requires, constexpr, consteval, constinit, noexcept,
                  nullptr, override, final, co_await, co_return, co_yield,
                  expected, unexpected, span, string_view, unique_ptr, shared_ptr},
    emph={maui, view, element, handler, layout, page, application, bindable_object,
          view_ref, weak_ref, observable, grid, stack_layout, content_page},
    emphstyle=\color{cppblue},
    inputencoding=utf8,
    % Не-ASCII знаци, срещащи се в коментарите на изходния текст. Без тези редове
    % \lstinputlisting отказва с „Invalid UTF-8 byte sequence“: listings не разбира
    % многобайтов вход, който не е в literate таблицата.
    literate={→}{{$\rightarrow$}}1 {×}{{$\times$}}1 {…}{{...}}1 {§}{{\S}}1
{а}{{\cyra}}1 {б}{{\cyrb}}1 {в}{{\cyrv}}1 {г}{{\cyrg}}1 {д}{{\cyrd}}1
             {е}{{\cyre}}1 {ж}{{\cyrzh}}1 {з}{{\cyrz}}1 {и}{{\cyri}}1 {й}{{\cyrishrt}}1
             {к}{{\cyrk}}1 {л}{{\cyrl}}1 {м}{{\cyrm}}1 {н}{{\cyrn}}1 {о}{{\cyro}}1
             {п}{{\cyrp}}1 {р}{{\cyrr}}1 {с}{{\cyrs}}1 {т}{{\cyrt}}1 {у}{{\cyru}}1
             {ф}{{\cyrf}}1 {х}{{\cyrh}}1 {ц}{{\cyrc}}1 {ч}{{\cyrch}}1 {ш}{{\cyrsh}}1
             {щ}{{\cyrshch}}1 {ъ}{{\cyrhrdsn}}1 {ь}{{\cyrsftsn}}1 {ю}{{\cyryu}}1 {я}{{\cyrya}}1
             {А}{{\CYRA}}1 {Б}{{\CYRB}}1 {В}{{\CYRV}}1 {Г}{{\CYRG}}1 {Д}{{\CYRD}}1
             {Е}{{\CYRE}}1 {Ж}{{\CYRZH}}1 {З}{{\CYRZ}}1 {И}{{\CYRI}}1 {Й}{{\CYRISHRT}}1
             {К}{{\CYRK}}1 {Л}{{\CYRL}}1 {М}{{\CYRM}}1 {Н}{{\CYRN}}1 {О}{{\CYRO}}1
             {П}{{\CYRP}}1 {Р}{{\CYRR}}1 {С}{{\CYRS}}1 {Т}{{\CYRT}}1 {У}{{\CYRU}}1
             {Ф}{{\CYRF}}1 {Х}{{\CYRH}}1 {Ц}{{\CYRC}}1 {Ч}{{\CYRCH}}1 {Ш}{{\CYRSH}}1
             {Щ}{{\CYRSHCH}}1 {Ъ}{{\CYRHRDSN}}1 {Ь}{{\CYRSFTSN}}1 {Ю}{{\CYRYU}}1 {Я}{{\CYRYA}}1,
}


\begin{document}

% ============================================================================
% Заглавна страница (А.5)
% ============================================================================
\begin{titlepage}
\begin{center}
    {\fontsize{18}{21.6}\selectfont\bfseries ТЕХНИЧЕСКИ УНИВЕРСИТЕТ \textendash{} СОФИЯ\par}
    \vspace{0.5cm}
    {\fontsize{16}{19.2}\selectfont\bfseries ФАКУЛТЕТ „КОМПЮТЪРНИ СИСТЕМИ И ТЕХНОЛОГИИ“\par}

    \vspace{0.3cm}
    \begin{figure}[H]
        \centering
        \includegraphics[width=0.30\linewidth]{images/tu-logo.png}
    \end{figure}
    \vspace{0.2cm}

    {\fontsize{20}{23}\selectfont\bfseries СИСТЕМА ЗА ПОДПОМАГАНЕ НА РЕШЕНИЯТА ПРИ ИЗБОР И ПОДРЕЖДАНЕ НА ПОРТФЕЙЛ ОТ ИТ ПРОЕКТИ ПРИ БЮДЖЕТНИ И РЕСУРСНИ ОГРАНИЧЕНИЯ\par}

    \vspace{0.6cm}

    % Формата на заглавието зависи от това какво възлага дисциплината. Учебният
    % план изброява и двете форми, „курсов проект“ и „курсова работа“, затова
    % редът по-долу може да се наложи да чете КУРСОВА РАБОТА. Да се потвърди
    % с преподавателя преди печат.
    {\fontsize{14}{16.8}\selectfont\bfseries КУРСОВ ПРОЕКТ\par}
    \vspace{0.3cm}
    {\fontsize{14}{16.8}\selectfont\bfseries по дисциплина „Информационни технологии за бизнес мениджмънт“\par}

    \vfill

    \begin{flushleft}
    {\fontsize{14}{16.8}\selectfont\bfseries Разработил:\par}
    Алекс Цветанов, фак. № \TODO{фак. номер}\par
    Специалност: Компютърно и софтуерно инженерство, ОКС „Магистър“\par

    \vspace{0.7cm}

    {\fontsize{14}{16.8}\selectfont\bfseries Преподавател:\par}
    Доц. д-р Иван Станков и Проф. д-р Димитър Вергиев\par
    \end{flushleft}

    \vfill

    % Годината е тази, в която проектът се предава. Не се попълва по предположение.
    {\fontsize{12}{14.4}\selectfont София, \TODO{година на предаване}\par}
\end{center}
\end{titlepage}

% ============================================================================
% Съдържание (А.3, А.6)
% ============================================================================
\newpage
\pagestyle{empty}
\setcounter{tocdepth}{2}
\tableofcontents
\newpage

\pagestyle{fancy}
\pagenumbering{arabic}

% ============================================================================
% УВОД, без номер
% ============================================================================
% УВОД, без номер. Проблемът и цената му вървят преди архитектурата.
\section*{УВОД}
\label{sec:introduction}
\addcontentsline{toc}{section}{УВОД}

Всяка организация с информационни технологии има повече идеи за проекти,
отколкото пари и хора да ги изпълни. Кой се финансира и в какъв ред решава
управлението. На практика решението се взема в електронна таблица, а таблицата
има три недостатъка, които не са естетически.

Първо, допусканията в нея не са записани никъде. Кой критерий колко тежи, защо
един проект е предпоставка за друг, откъде идва разчетната стойност: всичко живее
в главата на попълнилия клетките, и шест месеца по-късно никой, включително той,
не може да каже защо е така.

Второ, отхвърлянето идва без причина. Не се казва дали проектът е спрян от
бюджета, от липсата на инженери, от неизпълнена предпоставка или от по-ниската
оценка. Разликата е съществена: бюджетът се преразглежда идната година, а ниската
оценка не се.

Трето, и това е най-скъпото, таблицата смята с очаквани стойности, при които
портфейл, побиращ се в бюджета, изглежда безопасен. Настоящата разработка измери
какъв е той в действителност върху съставен сценарий: портфейл с очаквана стойност
1123 при бюджет 1150 има \textbf{38,7 процента} вероятност да остане в бюджета,
когато входните оценки се изтеглят от собствените си разпределения. Тоест около
три на пет шанс за преразход при решение, което изглежда с резерв. Таблицата с
очаквани стойности не може да покаже това число.

Проектът разработва система, наречена \textbf{Quorum}, която подпомага
решението, вместо да го замества. Тя приема бизнес процеса като насочен граф от
дейности със стойности, продължителности и зависимости, избира портфейла с
най-висока сумарна полза при зададените ограничения и преизчислява избора при
изменени допускания. Трите недостатъка определят трите изисквания: допусканията са
в текстов файл, който се версионира; всеки отхвърлен проект се връща с
ограничението, което го е спряло; рискът се смята върху разпределения.

\textbf{Целта} е класически методи на управленската наука, многокритериален анализ
и целочислена оптимизация, да бъдат вградени в работеща система така, че
препоръката да е проследима до входните допускания. \textbf{Обхватът} спира при
одобряването на портфейла: не са реализирани проследяване на изпълнението,
счетоводно отчитане, интеграция с ERP и уеб интерфейс.

\newpage

% ============================================================================
% I. Анализ на предметната област и литературен преглед
% ============================================================================
% I. Анализ на предметната област и литературен преглед.
% Това е разделът, който носи най-голяма тежест в този проект: темата е от
% управленската наука, а не от ежедневната инженерна практика на автора, затова
% всяко твърдение тук стъпва на цитиран източник.
\section{Анализ на предметната област и литературен преглед}
\label{sec:literature}

\subsection{Информационни технологии и подпомагане на управленските решения}

Разделението на управленските решения на структурирани, полуструктурирани и
неструктурирани е въведено от Gorry и Scott Morton~\cite{gorry1971framework} и
остава полезно и днес, защото определя каква роля може да поеме една
информационна система. Структурираното решение се автоматизира изцяло.
Неструктурираното остава изцяло на човека. Полуструктурираното, а изборът на
портфейл от проекти е точно такова, се поддава на изчислителна поддръжка само
частично: изчислима е следствената част, а преценката кои критерии колко тежат
остава управленска. Power~\cite{power2002dss} систематизира класовете системи за
подпомагане на решенията и подчертава същото: системата предоставя модел и
данни, но не поема отговорността за избора.

Оттук следва проектното изискване, което определя целия облик на разработваната
система. Тя не произвежда „правилния отговор“, а прави препоръката проследима:
за всяка избрана позиция трябва да е видимо кое ограничение я е допуснало и коя
оценка я е класирала. \TODO{допълни с втори източник за изискването за
проследимост на препоръката в системите за подпомагане на решенията}

\subsection{Управление и моделиране на бизнес процеси}

Съвременното управление на бизнес процеси се разглежда като жизнен цикъл от
идентификация, откриване, анализ, преработване, внедряване и наблюдение
\cite{dumas2018bpm}. За настоящия проект е съществена фазата на анализ, при
която процесът е вече описан формално и подлежи на количествена оценка по
време, стойност и ресурс. Нотацията BPMN 2.0~\cite{omg2013bpmn} е приетият
стандарт за такова описание и дава семантика на потока от дейности, на
разклоненията и на събитията.

Мрежовият модел на дейностите с продължителности и зависимости, върху който
стъпва изчисляването на най-ранните и най-късните срокове, е по-стар от BPMN и е
известен като метод на критичния път. Той дава на системата две неща наведнъж:
най-ранния възможен срок за завършване на един проект и запаса от време на всяка
негова дейност. Запасът е важен, защото дейност без запас превръща всяко
закъснение в закъснение на целия проект, а това е вход към оценката на риска.
\TODO{цитирай първоизточника за метода на критичния път, Kelley и Walker, 1959,
след като е проверен библиографски}

\subsection{Многокритериален анализ на решенията}

Проектите кандидати не се сравняват по един показател. Стратегическата стойност,
техническият риск, регулаторната необходимост и очакваната възвръщаемост са
несъизмерими и част от тях се измерват само по субективна скала. Методът на
аналитичната йерархия~\cite{saaty1980ahp, saaty1990decision} решава това чрез
попарни сравнения и извежда теглата на критериите от собствен вектор на
матрицата на сравненията, като предлага и мярка за съгласуваност на преценките.
Методът TOPSIS~\cite{hwang1981topsis} подхожда различно: класира алтернативите по
относителна близост до идеалното решение и отдалеченост от най-лошото, при
предварително зададени тегла. Triantaphyllou~\cite{triantaphyllou2000mcdm}
сравнява семейството методи и показва, че различните методи върху едни и същи
данни невинаги дават едно и също класиране. Именно това несъвпадение е причината
двата метода да бъдат реализирани и съпоставени в Раздел~\ref{sec:alternatives},
вместо единият да бъде избран предварително.

\subsection{Избор на портфейл и управление на инвестиционния риск}

Идеята, че портфейлът се оценява като цяло, а не като сбор от независимо добри
позиции, произхожда от Markowitz~\cite{markowitz1952portfolio}. Пренесена върху
портфейл от ИТ проекти, тя означава, че сумарната полза не е единственият
критерий: концентрацията на риск в няколко зависими проекта е свойство на
съвкупността. Стандартът за управление на портфейли~\cite{pmi2017portfolio}
описва управленската рамка около това решение, включително балансирането на
портфейла спрямо стратегическите цели.

От изчислителна страна задачата за избор на подмножество проекти с максимална
сумарна полза при ограничен бюджет е разновидност на задачата за раницата.
Класификацията на вариантите и точните алгоритми са изложени
в~\cite{martello1990knapsack} и~\cite{kellerer2004knapsack}. Задачата е
NP-трудна в общия си вид, но практическите размери в разглеждания случай са
малки, десетки проекти, а не милиони, което прави точното решаване напълно
приложимо. Добавянето на ограничение по капацитет за няколко периода и на
отношения на предхождане между проектите превръща формулировката в
многомерна задача за раницата с ограничения за предхождане.
\TODO{провери и цитирай източник за многомерния вариант с ограничения за
предхождане; не се позовавай на~\cite{kellerer2004knapsack} за твърдение, което
не е проверено в текста му}

\subsection{Изводи от прегледа}

Прегледът дава три опорни точки за проектирането. Първо, системата поддържа
решение, което остава на човека, затова изходът ѝ трябва да е обяснимо, а не
само оптимално множество. Второ, оценяването и подборът са две отделни задачи и
могат да се сменят независимо, което е аргумент за ясна граница между двата
слоя. Трето, изборът на метод за многокритериална оценка не е безразличен за
резултата, затова сравнението между два метода е част от проекта, а не
подробност.

\newpage

% ============================================================================
% II. Анализ на възможните решения и обосновка на избора
% ============================================================================
% II. Анализ на възможните решения и обосновка на избора.
% Три независими избора: метод за многокритериална оценка, метод за решаване на
% задачата за подбор, технологична основа. Всеки се решава отделно.
\section{Анализ на възможните решения и обосновка на избора}
\label{sec:alternatives}

\subsection{Метод за многокритериална оценка}

Разгледани са два метода. Методът на аналитичната йерархия~\cite{saaty1980ahp}
извежда теглата на критериите от попарни сравнения и дава на оценяващия мярка за
съгласуваност на собствените му преценки. Цената е обемът на въпросите: за $n$
критерия сравненията са $n(n-1)/2$, а при добавяне на нов критерий целият набор
се преразглежда. Методът TOPSIS~\cite{hwang1981topsis} изисква теглата да бъдат
дадени наготово и класира алтернативите по относителна близост до идеалната
точка, което го прави евтин при много алтернативи и стабилен при добавянето на
нова.

Изборът е двата метода да бъдат реализирани, а не единият. Причината е емпирична
и е посочена в прегледа: върху едни и същи данни различните методи могат да
класират различно~\cite{triantaphyllou2000mcdm}. Реализирането само на единия би
скрило този факт зад числото, което той произвежда. При двата реализирани метода
несъвпадението в класирането става наблюдаемо и се превръща в резултат на
проекта, а не в мълчаливо допускане. Аналитичната йерархия се използва за
извличане на теглата, а TOPSIS за класирането при тези тегла, с възможност и
двете да се сменят.

\subsection{Метод за решаване на задачата за подбор}

Задачата за избор на портфейл при бюджет, капацитет по периоди и отношения на
предхождане е формулирана като задача на целочисленото линейно програмиране.
Разгледани са три пътя за решаването ѝ.

\begin{enumerate}[topsep=0pt,itemsep=0pt,leftmargin=*]
\item \textbf{Лакома евристика} по съотношение полза към стойност. Евтина за
      реализация, но не дава гаранция за оптималност и не се справя с
      ограниченията за предхождане без допълнителни поправки.
\item \textbf{Собствен алгоритъм разклоняване и ограничаване} по схемата
      от~\cite{martello1990knapsack}. Дава точно решение и е напълно прозрачен,
      но за многомерния вариант с предхождане изисква съществено повече работа,
      отколкото проектът предвижда за този слой.
\item \textbf{Готов решател}, конкретно решателят за задачи с ограничения от
      пакета OR-Tools. Дава точно решение, поема многомерните ограничения без
      допълнителна работа и позволява същият модел да се преизползва при анализа
      на чувствителността.
\end{enumerate}

Избран е третият път. Основанието е разпределението на усилието: приносът на
проекта е в моделирането и в проследимостта на препоръката, не в реализацията на
решател. Лакомата евристика остава реализирана като контролна стойност, срещу
която се проверява дали решателят връща поне толкова добро решение. Това е и
единственият начин грешка в модела да бъде забелязана: решател, който връща
по-лош резултат от евристиката, не е бавен, а сгрешен.

\subsection{Технологична основа}

Ядрото на модела е реализирано на C++20. Изборът е обоснован от два фактора:
анализът на чувствителността изпълнява един и същ модел стотици пъти при
изменени входове, което поставя изисквания към времето за изпълнение, а пакетът
OR-Tools предоставя първокласен интерфейс на същия език. Алтернативата, Python с
интерфейса на същия пакет, би съкратила времето за писане, но би довела до втори
език в системата без друга полза.

Потребителският интерфейс е отделен от ядрото и общува с него чрез описание на
сценария в JSON. Границата е поставена там съзнателно: сценарият е файл, който
може да се версионира, да се сравнява и да се прилага към проекта като
доказателство какви са били допусканията. Интерфейсът е уеб базиран, за да не
въвежда зависимост от платформен инструментариум за изобразяване.
\TODO{потвърди окончателния избор на HTTP слой и го запиши тук; към момента на
писане решението не е фиксирано}

\newpage

% ============================================================================
% III. Проектиране на системата
% ============================================================================
% III. Проектиране на системата.
% Слоеве, модел на данните, формати. Диаграмите се добавят, когато слоевете са
% фиксирани. Празна фигура е по-лоша от липсваща.
\section{Проектиране на системата}
\label{sec:design}

\subsection{Слоеве и отговорности}

Системата е разделена на четири слоя с еднопосочна зависимост отгоре надолу.
Слоят на \term{модела} съдържа описанието на бизнес процеса и на проектите
кандидати и не зависи от нищо друго. Слоят на \term{оценяването} превръща
описанието в число за всеки проект чрез избрания метод за многокритериална
оценка. Слоят на \term{подбора} формулира и решава задачата за избор на
портфейл. Слоят на \term{представянето} приема сценарий, извиква останалите три
и връща резултата заедно с обяснението към него.

Границите между слоевете са поставени така, че смяната на метода за оценяване да
не засяга подбора, а смяната на алгоритъма за подбор да не засяга оценяването.
Това не е абстракция за всеки случай, а пряко следствие от
Раздел~\ref{sec:alternatives}: и двата слоя имат по повече от една реализация още
в първата версия. Потокът от данни между слоевете е показан на
Фиг.~\ref{fig:layers}.

\begin{figure}[H]
\centering
\begin{tikzpicture}[node distance=0.75cm and 1.1cm]
  \node[block] (scen) {Сценарий\\\code{scenarios/*.json}};
  \node[block, below=of scen] (model) {Слой на модела\\\code{model.cpp}, \code{cpm.cpp}\\граф, инварианти,\\критичен път, PERT};
  \node[block, below=of model] (score) {Слой на оценяването\\\code{mcdm.cpp}\\аналитична йерархия, TOPSIS};
  \node[block, below=of score] (select) {Слой на подбора\\\code{portfolio.cpp}};
  \node[accent, left=of select] (greedy) {Лакома евристика\\долна граница};
  \node[accent, right=of select] (exact) {Разклоняване и\\ограничаване\\линейна релаксация};
  \node[danger, below=of select] (check) {Проверка:\\точно $\ge$ евристика};
  \node[block, below=of check] (post) {Риск и чувствителност\\\code{risk.cpp}, \code{sensitivity.cpp}};
  \node[block, below=of post] (out) {Представяне\\\code{report.cpp}, \code{main.cpp}};

  \draw[line] (scen) -- (model);
  \draw[line] (model) -- (score);
  \draw[line] (score) -- (select);
  \draw[line] (select) -- (greedy);
  \draw[line] (select) -- (exact);
  \draw[line] (greedy) |- (check);
  \draw[line] (exact) |- (check);
  \draw[line] (check) -- (post);
  \draw[line] (post) -- (out);
\end{tikzpicture}
\caption{Слоеве и поток от данни. Стрелките сочат посоката на зависимостта.}
\label{fig:layers}
\end{figure}

Проверката преди последните два блока не е украса на диаграмата. Тя е
единственото място, където грешка в построяването на модела може да бъде уловена
автоматично: ако точният алгоритъм върне стойност под тази на допустимо решение,
намерено от евристиката, изпълнението спира със съобщение за грешка в модела,
вместо да отчете число.

\subsection{Модел на данните}

Бизнес процесът се представя като насочен ацикличен граф. Върховете са дейности
с продължителност, стойност и потребен ресурс. Ребрата са отношения на
предхождане. Проектът кандидат е свързано подмножество от дейности заедно с
оценките си по критериите. Портфейлът е подмножество от проекти.

Ацикличността е инвариант, който системата проверява при зареждане, а не
допускане. Цикъл в графа на предхождането прави най-ранният срок на завършване
недефиниран, а задачата за подбор неразрешима; откриването му при зареждането
дава на потребителя посочимо място в неговия вход, вместо съобщение за грешка от
алгоритъма. Проверката се прави на двете нива поотделно: между дейностите на
един проект и между самите проекти.

\begin{table}[H]
\centering
\caption{Обекти на модела и техните задължителни полета}
\label{tab:model}
\begin{tabular}{@{}L{3.0cm}L{4.4cm}L{4.6cm}@{}}
\toprule
\textbf{Обект} & \textbf{Задължителни полета} & \textbf{Инвариант} \\
\midrule
Дейност & идентификатор, продължителност, стойност & неотрицателни стойности, оптимистична $\le$ най-вероятна $\le$ песимистична \\
Зависимост & източник, цел & графът остава ацикличен, и двата края съществуват \\
Проект & идентификатор, дейности, оценки по критериите & непразно множество дейности, оценка по всеки критерий \\
Критерий & идентификатор, посока & сумата на теглата се нормира до единица \\
Ресурс & идентификатор, капацитет & неотрицателен капацитет \\
Ограничение & бюджет, капацитет по ресурс, предхождане, взаимно изключване & предхождането и изключването не се отнасят за една и съща двойка \\
\bottomrule
\end{tabular}
\end{table}

Съставът на модела е обобщен в Табл.~\ref{tab:model}. Посоката на критерия указва
дали по-голямата стойност е по-добра, което е необходимо и за нормализацията в
TOPSIS, и за обръщането на колоната при синтеза по аналитичната йерархия.

Три величини не се въвеждат, а се извеждат при зареждането, за да няма два
източника на едно и също число. Стойността на проекта е сумата от очакваните
стойности на дейностите му. Продължителността на проекта е дължината на неговия
критичен път. Потребният ресурс на проекта е върховото едновременно потребление
при разписание по най-ранни срокове, изчислено чрез обхождане на моментите на
започване и завършване. Последното е съзнателно консервативно: то допуска, че
всеки проект работи с върховото си потребление през целия период, в който е
активен, и следователно никога не подценява натоварването.

\subsection{Формат на сценария и на резултата}

Входът е JSON документ с фиксирана схема. Сценарият съдържа критериите, матрицата
на попарните сравнения, ресурсите с техните капацитети, бюджета и проектите с
дейностите и оценките им. Пълната схема е дадена в Приложение~А.

Изходът е текстов доклад с таблици, а не втори JSON документ. Причината е, че
изходът не се подава на друга програма: той се чете от човек и се прилага към
решението. За всеки отхвърлен проект се отпечатва причината за отхвърлянето, като
причините се търсят в определен ред: непокрито предхождане, взаимно изключване с
финансиран проект, недостиг на бюджет, недостиг на ресурс. Ако нито едно от тези
не важи, проектът е бил допустим и е отпаднал по стойност, и това се казва
изрично, вместо да се посочи ограничение, което не е било определящо.

\subsection{Анализ на риска}

Тройните оценки в сценария носят несигурност, която очакваната стойност изхвърля.
Анализът на риска я връща обратно, като извлича извадка от разпределението на
всяка продължителност и всяка стойност и преизчислява разписанието и сумарната
стойност за всяка извадка. Резултатът е разпределение на общата стойност на
портфейла и на срока на завършване на програмата, от което се отчитат средна
стойност, медиана и централен интервал.

Проектното решение тук е разпределението да бъде триъгълно върху трите точки.
То се определя изцяло от числата, които сценарият вече съдържа, няма
допълнителен параметър за форма, който да бъде предположен, и обратната му
функция е в затворен вид, което при стотици хиляди извадки има значение.
Разликата спрямо бета разпределението, което стои зад очакваната стойност на
PERT, е отчетена в Раздел~\ref{sec:implementation}.

\subsection{Анализ на чувствителността}

Анализът на чувствителността се проектира като повторно изпълнение на същия
модел върху систематично изменен вход, а не като отделен алгоритъм. Изменят се
теглата на критериите и бюджетът, всяко в зададен интервал и стъпка. Теглата се
изменят по едно наведнъж и след това се нормират наново, така че всяко изпълнение
да се отнася до едно-единствено изменено допускане; едновременното изменение би
дало резултат, който не може да бъде приписан на определена причина.

Резултатът е множество от портфейли, от които се извежда кои позиции остават
избрани при всички изменения и кои се появяват и изчезват. Първите са
устойчивата част от препоръката и това е основният изход, предназначен за
управленския читател. Освен обобщението по позиции се отпечатва и решетка с по
един стълб на изпълнение, защото делът сам по себе си казва колко често една
позиция е била финансирана, но не и при какви допускания.

\newpage

% ============================================================================
% IV. Програмна реализация
% ============================================================================
% IV. Програмна реализация.
% Пише се докато кодът се пише, не преди това. Твърденията тук трябва да сочат
% към съществуващ файл.
\section{Програмна реализация}
\label{sec:implementation}

\subsection{Обща структура на изходния текст}

Реализацията следва слоевете от Раздел~\ref{sec:design}. Всеки слой е отделна
преводна единица с публичен заглавен файл, а зависимостите между тях са
еднопосочни и се проверяват от системата за построяване. Външните зависимости са
две: пакетът за оптимизация и библиотеката за четене и записване на JSON.

\TODO{таблица на модулите: име, отговорност, публичен интерфейс; попълва се от
действителното дърво на изходния текст, не от плана}

\subsection{Слой на модела}

Слоят на модела съдържа представянето на графа и проверката на инвариантите му.
Проверката за ацикличност се извършва с топологично подреждане при зареждане на
сценария. Същото подреждане се преизползва при изчисляването на най-ранните
срокове, така че обходът на графа се прави веднъж, а не два пъти.

\TODO{коментар на изходния текст на топологичното подреждане и на изчисляването
на сроковете, с листинг}

\subsection{Слой на оценяването}

Двата метода за многокритериална оценка са зад общ интерфейс, който приема
матрицата проекти по критерии и връща вектор от оценки. Аналитичната йерархия
изчислява теглата от матрицата на попарните сравнения и отчита коефициента на
съгласуваност; при стойност над приетия праг оценяването се отказва вместо да
върне число, изведено от несъгласувани преценки. TOPSIS нормализира матрицата,
изчислява разстоянията до идеалната и до най-лошата точка и връща относителната
близост.

\TODO{запиши приетия праг за коефициента на съгласуваност и източника, от който е
взет; не го попълвай по памет}

\subsection{Слой на подбора}

Задачата се формулира с двоични променливи, по една на проект, целева функция
сумата от оценките и ограничения по бюджет, по капацитет за всеки период и по
предхождане. Лакомата евристика е реализирана в същия слой и се изпълнява
винаги, за да служи за долна граница при проверката на резултата от решателя.

\subsection{Слой на представянето}

Слоят на представянето приема сценарий, изпълнява останалите три слоя и връща
резултата. Обяснението към препоръката се сглобява тук: за всеки отхвърлен
проект се записва кое ограничение е било определящо.

\TODO{опиши крайните точки на интерфейса и формата на отговорите, след като са
реализирани}

\subsection{Тестове}

Проверката на реализацията е разделена на два вида. Модулните тестове покриват
отделните изчисления: топологичното подреждане, нормализацията в TOPSIS,
изчисляването на теглата. Тестовете върху цялата система изпълняват сценарии с
известно предварително решение и сравняват изхода с него.

\TODO{посочи достигнатото покритие и инструмента, с който е измерено}

\newpage

% ============================================================================
% V. Методика на експерименталната проверка
% ============================================================================
% V. Методика на експерименталната проверка.
% ВАЖНО: този раздел описва ПЛАНА. Никакви числа тук. Числата живеят в
% Раздел VI и се вписват едва след като са измерени.
\section{Методика на експерименталната проверка}
\label{sec:method}

\subsection{Какво се проверява}

Проверката отговаря на три въпроса, поставени независимо един от друг. Първо,
дава ли системата правилен отговор там, където правилният отговор е известен
предварително. Второ, съвпадат ли класиранията на двата метода за
многокритериална оценка върху едни и същи данни, и ако не, къде се разминават.
Трето, колко устойчива е препоръката спрямо изменение на входните допускания.

Първият въпрос е за коректност и има еднозначен отговор. Вторият и третият са
измервания, чиито резултати не се знаят предварително и не се предполагат тук.

\subsection{Набор от сценарии}

Проверката се извършва върху три групи сценарии. Първата съдържа малки сценарии,
за които оптималният портфейл е изчислен ръчно чрез изчерпателно изброяване;
тяхната роля е да уловят грешка в модела, а не в решателя. Втората съдържа
сценарии, породени по зададено правило, с контролиран брой проекти и плътност на
графа на зависимостите; те служат за наблюдение на поведението при нарастващ
размер. Третата съдържа поне един сценарий, съставен по действителна практика на
планиране.

\TODO{опиши произхода на сценария от практиката, или, ако такъв не бъде получен,
запиши изрично, че третата група отпада, и не я замествай с измислен пример}

\subsection{Процедура}

За всеки сценарий се изпълнява следната процедура. Сценарият се зарежда и
инвариантите му се проверяват. Оценките се изчисляват по всеки от двата метода.
За всеки вектор от оценки задачата за подбор се решава веднъж с лакомата
евристика и веднъж с решателя. Резултатите се записват заедно с изразходваното
време. Анализът на чувствителността се изпълнява отделно, като бюджетът се
изменя в зададен интервал със зададена стъпка, а теглата на критериите се
изменят по същия начин.

Условие за валидност: резултатът на решателя не може да бъде по-лош от този на
евристиката. Нарушение на това условие се третира като грешка в модела и спира
експеримента, а не се отчита като резултат.

\TODO{фиксирай интервалите и стъпките за изменение на бюджета и теглата, както и
броя повторения за измерването на времето}

\subsection{Форма на резултатите}

Резултатите се представят в три вида. Таблица с по един ред на сценарий, която
съдържа броя проекти, броя избрани позиции, сумарната оценка по двата метода и
времето за решаване. Таблица на разминаванията между двата метода за
многокритериална оценка, с по един ред на всяко разминаване в класирането.
Графика на устойчивостта, която показва при кой бюджет всяка позиция влиза в
препоръчания портфейл и при кой излиза от него.

\TODO{среда за измерване: машина, операционна система, компилатор, версия на
пакета за оптимизация; попълва се в деня на измерването}

\newpage

% ============================================================================
% VI. Експериментални резултати и анализ
% ============================================================================
% III. Експериментална проверка и резултати. Слети раздели V и VI.
% Всяко число тук е от изпълнение на quorum bench или quorum demo. Нищо не е
% пренесено по памет.
\section{Експериментална проверка и резултати}
\label{sec:results}

\subsection{Методика и среда}

Проверката отговаря на четири въпроса. Дава ли системата правилен отговор там,
където той е известен предварително. Съвпадат ли класиранията на двата метода за
оценяване и достига ли разминаването до портфейла. Каква е цената на точното
решаване спрямо евристиката. Колко устойчива е препоръката спрямо изменение на
допусканията. Използвани са две групи сценарии. Първата е породена от генератор с фиксирано
начално състояние, с контролиран брой проекти, две измерения на капацитета,
предхождане с вероятност една четвърт и няколко двойки взаимно изключване;
бюджетът е на четиридесет процента от исканото, за да е ограничението действащо.
Втората е един съставен сценарий, \code{scenarios/example.json}, с дванадесет
проекта, четири критерия, два ресурса и бюджет, покриващ около половината от
исканото. Трета група, сценарий по действителна практика в организация, беше
предвидена, но такъв не беше получен и групата отпада. Тя не се замества с
измислен пример: сценарий, съчинен така, че да изглежда действителен, не носи
повече доказателствена сила от породения по правило, а внушава, че носи.

Всички числа са получени на една и съща машина: AMD Ryzen 5 3600, 16 GiB,
Windows 11 Pro N, g++ 15.2.0 (MinGW-w64), \code{-std=c++20}, режим \code{Release},
CMake 4.3.2 и Ninja 1.13.2, без външни зависимости. Времето се измерва вътре в
програмата около самия алгоритъм, при едно и също нишково изпълнение.
Измерванията се изпълняват с \code{quorum bench}, а рискът и чувствителността с
\code{quorum demo}; и двете не приемат параметри, които биха променили отчетените
числа. Условие за валидност: точното решение не може да бъде по-лошо от
евристичното, а нарушение спира експеримента, вместо да се отчете като резултат.

\subsection{Препоръката и защо всеки проект е вътре или вън}

При бюджет 1150 и искане 2106 системата финансира седем от дванадесетте проекта
за 1123, тоест 97,7 процента от бюджета. Инженерният ресурс е изчерпан напълно,
26,0 от 26,0, докато операционният остава на 9,0 от 11,0. Търсенето обхожда 357
върха и доказва оптималност за 0,031 ms; лакомата евристика достига 97,55 процента
от оптималната стойност за 0,006 ms.

Критичният път минава през P03 и P04 и дава обща продължителност 34,7 периода.
Останалите пет финансирани проекта имат запас между 16,8 и 22,8 периода, тоест
забавяне в който и да е от тях не измества края на програмата.

Отхвърлянията идват с причина. И петте отхвърлени проекта се спират от бюджета:
остатъкът е 27, а цените им са между 62 и 325. Сред тях е P02, класиран първи и
по двата метода за оценяване. Предшественикът му P01 е финансиран, взаимно
изключване не го засяга, спира го единствено цената. Това е случаят, в който
подредбата по оценка и съставът на портфейла се разминават и системата дължи на
читателя причината, а не само списъка.

\subsection{Цената на решение, взето върху очаквани стойности}

При двадесет хиляди извадки върху избрания портфейл общата стойност има средна
1162 и стандартно отклонение 38,6, с деветдесетпроцентов централен интервал
$[1100; 1227]$. Срокът има средна 37,35 периода при стандартно отклонение 3,32 и
интервал $[32{,}02; 43{,}04]$. Оттук следва основният резултат на проекта:
портфейлът се побира в бюджета по очаквани стойности, 1123 срещу 1150, но
вероятността действителната сума да остане в бюджета е \textbf{38,7 процента},
тоест около три на пет шанс за преразход. Това е числото, което таблицата с
очаквани стойности не показва, и е причината анализът на риска да е част от
системата, а не допълнение към нея.

Разминаването между 1123 и 1162 е 3,5 процента и е следствие от съзнателно
избраното разпределение, а не от грешка: подборът работи с очаквани стойности по
PERT, при които модата получава тегло четири, а извадката се тегли от триъгълно
разпределение със средна, равна на средно аритметично на трите точки, а оценките
в сценария са несиметрични надясно. Средният срок надхвърля детерминираните 34,7
по същата причина, плюс ефекта на сливането на пътищата.

Броят извадки по подразбиране не е избран по усещане: оценката за p95 е повторена
с десет начални състояния на генератора и размахът пада от 7,99 при хиляда
извадки до 1,20 при сто хиляди.

\subsection{Устойчивост на препоръката}

При изменение на теглата, шестнадесет изпълнения, пет позиции остават в портфейла
винаги: P03, P04, P07, P09 и P11. P01 и P08 излизат по веднъж, P06 и P10 влизат
само по веднъж, а P02, P05 и P12 не влизат никога. При изменение на бюджета от
шестдесет до сто и четиридесет процента, девет изпълнения, картината е
по-подвижна и е дадена на Фиг.~\ref{fig:stability}.

\begin{figure}[H]
\centering
\begin{tikzpicture}[x=1.08cm, y=0.5cm]
  \foreach \r/\name in {1/P01,2/P02,3/P03,4/P04,5/P05,6/P06,7/P07,8/P08,9/P09,10/P10,11/P11,12/P12}
    \node[anchor=east, font=\small] at ({0.45}, {-\r}) {\name};
  \foreach \c/\f in {1/{0,6},2/{0,7},3/{0,8},4/{0,9},5/{1,0},6/{1,1},7/{1,2},8/{1,3},9/{1,4}}
    \node[anchor=south, font=\small] at ({\c}, {-0.6}) {\f};
  \foreach \r in {1,...,12}
    \foreach \c in {1,...,9}
      \draw[gray!35, line width=0.3pt] ({\c-0.42}, {-\r-0.2}) rectangle ({\c+0.42}, {-\r+0.2});
  \foreach \c/\r in {1/1,2/1,3/1,5/1,6/1,7/1,8/1,9/1,
                     2/2,3/2,
                     4/3,5/3,6/3,7/3,8/3,9/3,
                     4/4,5/4,6/4,7/4,8/4,9/4,
                     6/6,7/6,8/6,9/6,
                     1/7,2/7,3/7,4/7,5/7,6/7,7/7,8/7,9/7,
                     1/8,2/8,3/8,4/8,5/8,6/8,7/8,8/8,9/8,
                     1/9,2/9,3/9,4/9,5/9,6/9,7/9,8/9,9/9,
                     1/10,4/10,
                     3/11,4/11,5/11}
    \fill[blue!45] ({\c-0.42}, {-\r-0.2}) rectangle ({\c+0.42}, {-\r+0.2});
  \node[anchor=north, font=\small] at ({5}, {-12.8}) {бюджет като дял от зададения};
\end{tikzpicture}
\caption{Позиции в избрания портфейл при изменение на бюджета. Запълнена клетка означава финансиран проект.}
\label{fig:stability}
\end{figure}

Три позиции, P07, P08 и P09, са в портфейла при всеки от деветте бюджета. Две,
P05 и P12, не влизат при нито един, включително при сто и четиридесет процента:
те губят по оценка, а не само по бюджет. Между тях са седемте подвижни позиции.

Фигурата показва и нещо, което делът в проценти би скрил: присъствието на проект
не е монотонно по бюджета. P01 е финансиран при шестдесет, седемдесет и осемдесет
процента, отпада при деветдесет и се връща при сто; P10 е финансиран при
шестдесет и при деветдесет, но не между тях. Причината е, че при по-широк бюджет
става възможна двойката P03 и P04, която измества други позиции: сумарната
стойност расте монотонно, но съставът не може. Оттук следва предупреждението за
управленския читател. Изречение като „P01 е препоръчан в осемдесет и девет
процента от случаите“ е вярно, но подвеждащо, защото внушава интервал, който не
съществува. Прочитът изисква решетката, а не само процента, затова системата
отпечатва и двете.

\subsection{Двата метода за оценяване дават един и същ портфейл}

Теглата, изведени от матрицата на попарните сравнения, са $0{,}1891$ за
стратегическото съответствие, $0{,}3509$ за възвръщаемостта, $0{,}3509$ за
регулаторната необходимост и $0{,}1091$ за риска, при коефициент на съгласуваност
$0{,}0036$, далеч под прага $0{,}10$. При тези тегла двата метода класират
дванадесетте проекта различно на шест позиции: P09, P06 и P04 излизат по-напред
по TOPSIS съответно с една, две и една позиция, а P10, P12 и P03 по-назад с две,
една и една. Коефициентът на рангова корелация на Spearman е $0{,}9580$; всички
разминавания са между пето и десето място.

Отражението върху решението е нулево: изпълнен поотделно с двата вектора от
оценки, подборът връща един и същ портфейл, P01, P03, P04, P07, P08, P09 и P11.
При този сценарий изборът между двата метода не е решаващ, и това е по-полезно за
управленския читател от която и да е от двете подредби поотделно. Обратното, че
методът никога не е от значение, не следва оттук: проверено е върху един сценарий.

\subsection{Коректност и цена на точното решаване}

Стойността от разклоняването и ограничаването е сравнена с изчерпателно изброяване
на всички подмножества при 8, 12, 16 и 18 проекта, по двадесет задачи на размер.
Съвпадението е пълно при всичките осемдесет. Средното време на задача расте от
0,0039 до 0,1069 ms, докато изброяването расте от 0,021 до 32,353 ms; при двадесет
и пет проекта еталонът би отнел часове, което е и горната граница на проверката.

\begin{table}[H]
\centering
\caption{Разклоняване и ограничаване срещу лакома евристика, по пет задачи на размер}
\label{tab:exact-greedy}
\small
\begin{tabular}{@{}cccccc@{}}
\toprule
\textbf{Проекти} & \textbf{Точно} & \textbf{Лакомо} & \textbf{Дял, \%} & \textbf{Върхове} & \textbf{Точно, ms} \\
\midrule
10 & 0,5753 & 0,5197 & 89,87 & 74     & 0,008 \\
14 & 0,7875 & 0,6940 & 87,84 & 225    & 0,018 \\
18 & 1,1170 & 0,9141 & 81,11 & 920    & 0,077 \\
22 & 1,2272 & 1,0708 & 87,43 & 3440   & 0,237 \\
26 & 1,5857 & 1,2652 & 79,01 & 7937   & 0,753 \\
30 & 1,8450 & 1,6092 & 86,58 & 18439  & 1,710 \\
34 & 2,2365 & 1,8923 & 84,75 & 54127  & 5,530 \\
38 & 2,4237 & 2,1413 & 88,53 & 73295  & 7,862 \\
\bottomrule
\end{tabular}
\end{table}

Евристиката губи между десет и двадесет и един процента от оптималната стойност,
без тенденция с размера, което е характерно за лаком избор при действащи
ограничения за предхождане. Цената на точността е малка: при тридесет и осем
проекта търсенето отнема под осем милисекунди, тоест анализът на чувствителността
остава под секунда и при задачи, много по-големи от съставния сценарий.
Отношението между двата алгоритъма обаче расте от около единица при десет проекта
до около триста и петдесет при тридесет и осем, тоест границата на точното
решаване е в порядъка на десетки проекти.

\newpage

% ============================================================================
% ЗАКЛЮЧЕНИЕ, без номер
% ============================================================================
% ЗАКЛЮЧЕНИЕ, без номер.
% Обобщава характеристиките на решението, предимствата и недостатъците му,
% резултатите и посоките за по-нататъшна работа.
\section*{ЗАКЛЮЧЕНИЕ}
\label{sec:conclusion}
\addcontentsline{toc}{section}{ЗАКЛЮЧЕНИЕ}

Проектът разработва система за подпомагане на решенията при избор и подреждане
на портфейл от ИТ проекти. Бизнес процесът се описва като насочен ацикличен граф
от дейности със стойности, продължителности и зависимости. Проектите кандидати
се оценяват по няколко критерия чрез два независими метода, а изборът на
портфейл се формулира като задача на целочисленото линейно програмиране с
ограничения по бюджет, по капацитет и по предхождане. Към препоръката се прилага
анализ на чувствителността, който отделя устойчивата ѝ част от частта, която се
мени заедно с допусканията.

\textbf{Предимства на решението.} Допусканията са в текстов файл, който се
версионира и сравнява, така че препоръката е възпроизводима. Оценяването и
подборът са отделени, което позволява методът за оценяване да бъде сменен, без
да се пипа подборът. Лакомата евристика, изпълнявана успоредно с решателя, дава
евтина проверка, че моделът не е сгрешен.

\textbf{Недостатъци и ограничения.} Системата приема оценките по критериите за
дадени и не проверява как са получени; качеството на препоръката не може да
надхвърли качеството на този вход. Формулировката е статична и не отчита
възможността даден проект да бъде отложен вместо отхвърлен. Обхватът спира при
одобряването на портфейла и не покрива проследяването на изпълнението.

\TODO{обобщение на действителните експериментални резултати; пише се след
Раздел~\ref{sec:results} и не го изпреварва}

\textbf{Посоки за по-нататъшна работа.} Първата е добавянето на времево
измерение, при което решението не е дали един проект да бъде изпълнен, а в кой
период да започне. Втората е вграждането на несигурността в оценките, чрез
интервали вместо точкови стойности. Третата е свързване с работеща система за
планиране на ресурсите на предприятието, така че разчетните стойности и
наличният капацитет да се четат оттам, а не да се въвеждат ръчно.

\newpage

% ============================================================================
% Списък на използваните литературни източници
% ============================================================================
\phantomsection
\addcontentsline{toc}{section}{ЛИТЕРАТУРА}
\bibliographystyle{ieeetr}
\bibliography{references}
\newpage

% ============================================================================
% Приложения
% ============================================================================
% ПРИЛОЖЕНИЕ, без номер на раздел.
% Изходният текст НЕ се вмъква тук. Хранилището е публично и една таблица с
% указател към файловете е по-полезна от двадесет страници преписан код.
\section*{ПРИЛОЖЕНИЕ}
\label{sec:appendix}
\addcontentsline{toc}{section}{ПРИЛОЖЕНИЕ}

Целият изходен текст, вложеният сценарий, тестовете, уеб клиентът и пълният
изход на командите \code{quorum demo} и \code{quorum bench} са в публичното
хранилище \url{https://github.com/Alex-Tsvetanov/quorum}. Построяване и
изпълнение, без външни зависимости, само с компилатор за C++20 и CMake:
\code{cmake -S . -B build -DCMAKE\_BUILD\_TYPE=Release}, след това
\code{cmake --build build} и \code{./build/quorum demo}. Уеб страницата се вдига с
\code{./build/quorum serve}. Указателят кое изчисление в кой файл живее е в
Табл.~\ref{tab:sources}; броят редове на ядрото е от \code{wc -l} към състоянието,
от което са получени числата в Раздел~\ref{sec:results}, а редовете на
представянето са към състоянието с уеб клиента.

\begin{table}[H]
\centering
\caption{Изходен текст: къде живее всяко изчисление}
\label{tab:sources}
\small
\begin{tabular}{@{}L{3.2cm}L{7.4cm}r@{}}
\toprule
\textbf{Файл в \code{src/}} & \textbf{Съдържание} & \textbf{Редове} \\
\midrule
\code{json.cpp} & четец и записвач на JSON с коментари & 368 \\
\code{model.cpp} & сценарий, инварианти, изведени величини & 271 \\
\code{cpm.cpp} & топологично подреждане, критичен път, PERT & 245 \\
\code{mcdm.cpp} & тегла по AHP, съгласуваност, TOPSIS & 244 \\
\code{portfolio.cpp} & разклоняване и ограничаване, лаком алгоритъм & 345 \\
\code{risk.cpp} & Монте Карло, триъгълно разпределение & 162 \\
\code{sensitivity.cpp} & сканиране по тегла и по бюджет & 132 \\
\code{report.cpp} & таблици и обяснения на отхвърлянето & 215 \\
\code{analyse.cpp} & общият конвейер към доклада и към JSON & 213 \\
\code{http.cpp}, \code{net.cpp}, \code{web.cpp} & HTTP/1.1, сокети, маршрутизация & 580 \\
\code{bench.cpp}, \code{main.cpp} & измерванията, командният ред и \code{serve} & 407 \\
\midrule
\code{include/quorum/} & дванадесет публични заглавни файла & 803 \\
\code{tests/} & 83 случая в девет групи и рамката за тестване & 1432 \\
\code{web/} & страницата: HTML, CSS и JavaScript без пакетна система & 538 \\
\code{scenarios/} & съставният сценарий с дванадесет проекта & 206 \\
\bottomrule
\end{tabular}
\end{table}

\textbf{Схема на сценария.} Сценарият е един JSON документ. Задължителни са
\code{budget}, \code{criteria} и \code{projects}. Критерият носи \code{id} и по
избор \code{direction}, \code{max} или \code{min}; теглата се задават пряко чрез
\code{weight} или се извеждат от квадратната положителна реципрочна матрица
\code{comparisons}. Ресурсът носи \code{id} и \code{capacity}. Проектът носи
\code{id}, оценка по всеки критерий в \code{scores}, по избор \code{requires} и
\code{excludes}, и непразен списък \code{activities}. Дейността носи \code{id},
\code{depends\_on} и \code{demand} по ресурс. Продължителността и стойността се
задават или с едно число, \code{duration} и \code{cost}, или с тройна оценка,
\code{optimistic}, \code{likely} и \code{pessimistic}, със същите представки за
стойността; едното число е случаят, в който трите точки съвпадат, затова за него
няма отделен път в кода.


\end{document}
